% !TEX root  = identitti_aufsatz.tex

\documentclass[12pt,a4paper]{article}

\usepackage[a4paper,margin=1in]{geometry}
\usepackage[T1]{fontenc}
\usepackage[utf8]{inputenc}
\usepackage{newtxtext,newtxmath}
\usepackage[hidelinkcs]{hyperref}
\usepackage{setspace}
\usepackage[ngerman]{babel}
\usepackage{csquotes}
\usepackage[backend=biber,style=verbose-note,ibidtracker=true]{biblatex}
\addbibresource{references.bib}

% Replace these values with your own content.
\newcommand{\CoverCourse}{Deutsch 328}
\newcommand{\CoverInstructor}{Herr Lukas Valtin}
\newcommand{\CoverDate}{den 1. Juni 2026}
\newcommand{\CoverDepartment}{Department of Modern Foreign Languages}
\newcommand{\CoverStudent}{Barnard, Quin Charl \#27902277}
\newcommand{\CoverEmail}{27902277@sun.ac.za}
\renewcommand{\contentsname}{Inhaltsverzeichnis}

\begin{document}

\begin{titlepage}
\thispagestyle{empty}

\vspace*{3.55in}

{\raggedleft
    {\fontsize{31}{35}\selectfont *TITLE HERE*\par}
\par}

\vspace{0.30in}

{\raggedleft
    {\itshape\fontsize{18}{22}\selectfont Eine Besprechung zum Roman \textit{„Identitti“}\par}
\par}

\vspace{0.45in}

\begin{raggedleft}
   {\fontsize{18}{21}\selectfont \CoverCourse\par}
    \vspace{0.03in}
    {\fontsize{11}{14}\selectfont
      \CoverInstructor\par
      \CoverDate\par
      \CoverDepartment\par}
\end{raggedleft}

\vspace{0.55in}

\begin{raggedleft}
    {\fontsize{15}{18}\selectfont \CoverStudent\par}
    \vspace{0.04in}
    {\fontsize{10}{12}\selectfont \CoverEmail\par}
\end{raggedleft}
\end{titlepage}

\setcounter{page}{1}
\tableofcontents
\clearpage

\onehalfspacing

\section{Einleitung}
{\centering
\itshape
\footnotesize
„[...] bloß weil ich keinen neuen Pass, keine neue Kategorie bekomme, sagt ihr, es gäbe das gar nicht, was ich getab habe. Aber es gibt mich und Menschen wie mich!“ \footcite[275]{sanyal2021identitti}
\par}
\noindent In Mithu Sanyals Roman \textit{„Identitti“}, der im Jahr 2021 veröffentlicht worden ist, wird der Begriff der Identität hinterfragt. Mannigfaltige Perspektive werden von Sanyal gegeneinander gestellt, um darzustellen, dass eine weitgehend inklusive, aber immer strukturierte Interpretation der Identität in der modernen Welt notwendig ist. Den Identitätsentwurf unterstützt die Struktur des Romans, da verschiedene Textsorten beisammen geführt werden um ein Werk mit autofiktiven Elementen zu kreiern. Der Roman stellt einen wichtigen Beitrag zur zeitgenössischen Diskussion der Frage der „Rasse.“

\newpage 

\section{Hauptteil}
Um den Roman erfolgreich deuten zu können, muss der Leser die formalen Aspekte dessen verstehen.

\subsection{Example Subsection}
Continue writing the body of the document after the table of contents.

\section{Conclusion}
Summarize your document here.

\end{document}
